\section{Experimental Analysis}

\subsection{Modeling of Primary Symbolic Signals}
The experimental procedure began with the mathematical definition of two distinct signals, $x(t)$ and $y(t)$. The first signal, $x(t)$, was constructed as a unit ramp function $t$ bounded by the interval $[-2, 2]$, while the second signal, $y(t)$, was defined as a reflected and scaled ramp $-2t$ within the narrower interval $[-1, 1]$. By utilizing Heaviside step functions to define these temporal windows, the signals were established with infinite mathematical precision, avoiding the discretization errors inherent in numerical methods. The initial plots in Figure 1 confirmed the accuracy of these definitions, showing $x(t)$ as the stationary reference and $y(t)$ as the waveform to be shifted for the correlation process.

\begin{figure}[h]
    \centering
    \includegraphics[width=\columnwidth]{figure1.png}
    \caption{Primary Signal}
    \label{fig:primary_signal}
\end{figure}

\subsection{Evaluation of Cross-Correlation Boundary Cases}
The cross-correlation $R_{xy}(\tau)$ was analyzed by evaluating the sliding integration of $x(t)$ and $y(t - \tau)$ across five specific boundary conditions. These cases represent the stages of signal interaction: 
\begin{itemize}
    \item \textbf{Case 1} ($\tau < -3$) and \textbf{Case 5} ($\tau \geq 3$): total separation with zero overlap.
    \item \textbf{Case 2} ($-3 \leq \tau < -1$) and \textbf{Case 4} ($1 \leq \tau < 3$): partial entry and exit of the signals.
    \item \textbf{Case 3} ($-1 \leq \tau < 1$): full overlap where $y(t)$ is entirely contained within the bounds of $x(t)$.
\end{itemize}
Following the identification of these cases, the integration of the product of the signals was performed for each interval. The resulting total cross-correlation function, rendered in Figure 2, combined these piecewise results into a single symmetrical waveform that effectively quantifies the degree of similarity between the two signals as a function of the time-lag $\tau$.

\begin{figure}[h]
    \centering
    \includegraphics[width=\columnwidth]{figure2.png}
    \caption{Cross-Correlation Cases and Total}
    \label{fig:cross_correlation}
\end{figure}

\subsection{LTI System Response via Convolution}
The final stage of the analysis focused on the behavior of a Linear Time-Invariant (LTI) system subjected to a specific input pulse. The input signal $x_{in}(t)$ was defined as a pulse-like waveform $u(t) - 2u(t-1) + u(t-2)$, and the system was characterized by an exponential impulse response $h(t) = e^{-10t}u(t)$. These signals were first plotted together in Figure 3 to visualize the system's initial state. The convolution was then evaluated through two temporal cases: \textbf{Case 1} ($t < 0$), where no interaction occurs, and \textbf{Case 2} ($t \geq 0$), where the input pulse is processed by the system's exponential decay. The final output response $y_{out}(t)$, as shown in Figure 3, exhibited a sharp rise followed by a rapid nonlinear decay, accurately reflecting the system's time constant and its causal response to the input stimulus.

\begin{figure}[h]
    \centering
    \includegraphics[width=\columnwidth]{figure3.png}
    \caption{LTI System Cases of Convolution and Output}
    \label{fig:lti_system}
\end{figure}